% Options for packages loaded elsewhere
\PassOptionsToPackage{unicode}{hyperref}
\PassOptionsToPackage{hyphens}{url}
%
\documentclass[
]{article}
\usepackage{amsmath,amssymb}
\usepackage{iftex}
\ifPDFTeX
  \usepackage[T1]{fontenc}
  \usepackage[utf8]{inputenc}
  \usepackage{textcomp} % provide euro and other symbols
\else % if luatex or xetex
  \usepackage{unicode-math} % this also loads fontspec
  \defaultfontfeatures{Scale=MatchLowercase}
  \defaultfontfeatures[\rmfamily]{Ligatures=TeX,Scale=1}
\fi
\usepackage{lmodern}
\ifPDFTeX\else
  % xetex/luatex font selection
\fi
% Use upquote if available, for straight quotes in verbatim environments
\IfFileExists{upquote.sty}{\usepackage{upquote}}{}
\IfFileExists{microtype.sty}{% use microtype if available
  \usepackage[]{microtype}
  \UseMicrotypeSet[protrusion]{basicmath} % disable protrusion for tt fonts
}{}
\makeatletter
\@ifundefined{KOMAClassName}{% if non-KOMA class
  \IfFileExists{parskip.sty}{%
    \usepackage{parskip}
  }{% else
    \setlength{\parindent}{0pt}
    \setlength{\parskip}{6pt plus 2pt minus 1pt}}
}{% if KOMA class
  \KOMAoptions{parskip=half}}
\makeatother
\usepackage{xcolor}
\usepackage{longtable,booktabs,array}
\usepackage{calc} % for calculating minipage widths
% Correct order of tables after \paragraph or \subparagraph
\usepackage{etoolbox}
\makeatletter
\patchcmd\longtable{\par}{\if@noskipsec\mbox{}\fi\par}{}{}
\makeatother
% Allow footnotes in longtable head/foot
\IfFileExists{footnotehyper.sty}{\usepackage{footnotehyper}}{\usepackage{footnote}}
\makesavenoteenv{longtable}
\setlength{\emergencystretch}{3em} % prevent overfull lines
\providecommand{\tightlist}{%
  \setlength{\itemsep}{0pt}\setlength{\parskip}{0pt}}
\setcounter{secnumdepth}{-\maxdimen} % remove section numbering
\ifLuaTeX
  \usepackage{selnolig}  % disable illegal ligatures
\fi
\IfFileExists{bookmark.sty}{\usepackage{bookmark}}{\usepackage{hyperref}}
\IfFileExists{xurl.sty}{\usepackage{xurl}}{} % add URL line breaks if available
\urlstyle{same}
\hypersetup{
  hidelinks,
  pdfcreator={LaTeX via pandoc}}

\author{}
\date{}

\begin{document}

\hypertarget{pear-holepunch-p2p-resolver-threat-model}{%
\section{Pear / Holepunch P2P Resolver Threat
Model}\label{pear-holepunch-p2p-resolver-threat-model}}

This document outlines the security assumptions, privacy rules, and
threat vectors for the SatsPath PearResolver integration.

\hypertarget{overview}{%
\subsection{1. Overview}\label{overview}}

The PearResolver allows SatsPath to resolve aliases (e.g.,
\texttt{alice@example.com}) to \texttt{SignedPaymentProfile}s over the
Holepunch/Hyperswarm peer-to-peer network. This bypasses centralized
servers and DNS, relying on distributed hash tables (DHT) to locate
peers announcing a specific topic.

\hypertarget{privacy-rules-p2p-03}{%
\subsection{2. Privacy Rules (P2P-03)}\label{privacy-rules-p2p-03}}

To prevent passive scraping of the Hyperswarm DHT, the following privacy
rules are strictly enforced:

\begin{enumerate}
\def\labelenumi{\arabic{enumi}.}
\tightlist
\item
  \textbf{No Plaintext Aliases in Topics}: The Hyperswarm topic used for
  discovery MUST NOT be the plaintext alias. It must be a cryptographic
  hash of the alias (e.g., \texttt{SHA256(alias)}).
\item
  \textbf{No Private Keys in Pear Node}: The \texttt{satspath-pear}
  daemon or script only handles public \texttt{SignedPaymentProfile}
  JSON objects. It never loads, generates, or transmits the user's
  \texttt{identity\_secret\_key} or any wallet secrets.
\item
  \textbf{Fail-Closed Resolution}: If a peer returns malformed data or a
  profile belonging to a different alias, it is dropped silently.
\end{enumerate}

\hypertarget{threat-vectors-and-mitigations}{%
\subsection{3. Threat Vectors and
Mitigations}\label{threat-vectors-and-mitigations}}

\begin{longtable}[]{@{}
  >{\raggedright\arraybackslash}p{(\columnwidth - 4\tabcolsep) * \real{0.1476}}
  >{\raggedright\arraybackslash}p{(\columnwidth - 4\tabcolsep) * \real{0.8118}}
  >{\raggedright\arraybackslash}p{(\columnwidth - 4\tabcolsep) * \real{0.0406}}@{}}
\toprule\noalign{}
\begin{minipage}[b]{\linewidth}\raggedright
Threat
\end{minipage} & \begin{minipage}[b]{\linewidth}\raggedright
Mitigation
\end{minipage} & \begin{minipage}[b]{\linewidth}\raggedright
Layer
\end{minipage} \\
\midrule\noalign{}
\endhead
\bottomrule\noalign{}
\endlastfoot
\textbf{Passive DHT Scraping} & Topics are \texttt{SHA256(alias)},
making it impossible to harvest a list of registered emails/aliases by
simply listening to the DHT. & Pear / JS \\
\textbf{Profile Spoofing / Man-in-the-Middle} & The Rust resolver
verifies the \texttt{secp256k1} signature of every profile returned by a
peer against the declared \texttt{identity\_pubkey}. A malicious peer
cannot forge a signature for an identity they do not control. & Rust
Core \\
\textbf{Replay Attacks} & The \texttt{SignedPaymentProfile} includes an
\texttt{expires\_at} timestamp. The Rust resolver rejects expired
profiles. & Rust Core \\
\textbf{Denial of Service (DHT Spam)} & An attacker could spam the topic
with junk data. The \texttt{satspath-pear} resolver will disconnect from
peers sending non-JSON or invalid schema data. The Rust resolver handles
validation quickly and drops invalid signatures. & Pear + Rust \\
\textbf{Alias Collision (Intentional)} & Two users cannot easily claim
the same alias because the sender verifies the profile signature against
the expected identity, or trust is established via out-of-band exchange
(TOFU or verified channels). & Rust Core \\
\end{longtable}

\hypertarget{trust-model}{%
\subsection{4. Trust Model}\label{trust-model}}

\begin{itemize}
\tightlist
\item
  \textbf{Hyperswarm Network}: Untrusted. Any node can join and announce
  any topic.
\item
  \textbf{satspath-pear Process}: Trusted (runs locally on the user's
  machine).
\item
  \textbf{Profile Data}: Untrusted until the Rust Core validates the
  \texttt{secp256k1} signature and expiration.
\end{itemize}

\end{document}
